\documentclass[10pt,a4paper]{amsart}
\usepackage[margin=19mm]{geometry}
\usepackage{amsmath,amssymb,mathtools}
\usepackage[scr=rsfs,cal=euler]{mathalfa}
\usepackage{xfrac}
\usepackage[unicode,hidelinks,bookmarks=false]{hyperref}
\newcommand{\ie}{i.e.\ }
\newcommand{\cf}{cf.\ }
\newcommand{\resp}{respectively\ }
\newcommand{\Subsection}{\subsection}
\providecommand{\nobreakdash}{\nobreak\hbox{-}}
\setlength{\emergencystretch}{2em}
\multlinegap0pt
\usepackage{amssymb}
\usepackage[shortlabels]{enumitem}
\usepackage{float}
\newtheorem{assumption}{Assumption}
\newtheorem{proposition}{Proposition}
\title{Constructal Evolution as a Nonsmooth Dynamical System: Stability and Selection of Flow Architectures}
\author{Pascal Stiefenhofer}
\date{Source: arXiv:2603.06705v1 (CC BY 4.0)}
\begin{document}
\maketitle

\noindent\emph{Source and licence.} Constructal Evolution as a Nonsmooth Dynamical System: Stability and Selection of Flow Architectures. Version arXiv:2603.06705v1 (CC BY 4.0), distributed by arXiv under the Creative Commons Attribution 4.0 International licence, http://creativecommons.org/licenses/by/4.0/, as recorded in the arXiv licence metadata for this exact version and shown on the abstract page https://arxiv.org/abs/2603.06705v1. Full licence notice: \texttt{licenses/arxiv-2603.06705-CC-BY-4.0.txt}.\par\smallskip

\noindent\textit{Editorial scope.} Counted unit: the article's proposition prop:dissipation\_sufficient (source0.tex lines 818--857). The article's complete original proof of it is delivered with it (source0.tex lines 859--976, one \texttt{proof} environment of 118 lines). No mathematics is written, repaired or re-derived: the statement and the proof are the article's own lines, in the article's own notation, and no retained line is substituted or re-flowed.  Retained uncounted as the source-local context the proof needs: lines 275--289.  Nothing was omitted from any retained span, so the package carries no omission record.  No retained line contains a citation, so the wrapper carries no bibliography and no bibliography entry.  No retained line contains a figure environment, an image-inclusion command or drawing code, so no image is delivered and the package declares no asset dependency.  Editorial formatting only, in addition: an AMS \texttt{amsart} preamble that restates the article's own declarations and macros used by the retained lines, namely the environments \texttt{assumption}, \texttt{proposition} on separate counters, together with the standard packages those lines need.  The retained environments print in this document as Assumption 1, Proposition 1 in order of appearance, whereas the article prints the same items with the numbering of its own full document.  The complete native source of the article as distributed in the arXiv e-print for this exact version is bundled unchanged as \texttt{sources/195890/source0.tex}.
\begin{assumption}[Constructal dissipation]
\label{ass:dissipation}
There exist $\alpha>0$ and a continuous function
\[
\Psi:K\to\mathbb{R}_{\ge0}
\]
such that
\begin{equation}
\sup_{v\in F(x)}
\mathcal{R}^\circ(x;v)
\le -\alpha\,\Psi(x),
\qquad x\in K .
\label{eq:dissipation_condition}
\end{equation}
\end{assumption}

\begin{proposition}[Projected subgradient sufficiency]
\label{prop:dissipation_sufficient}
Let $K\subset\mathbb{R}^n$ be nonempty, closed, and convex, and let
$\mathcal R:K\to\mathbb{R}$ be locally Lipschitz.
Define the set-valued map
\begin{equation}
F(x)
:=
\Pi_{T_K(x)}\!\big(-\partial_C \mathcal R(x)\big),
\qquad x\in K,
\label{eq:projected_subgradient}
\end{equation}
where $T_K(x)$ denotes the Bouligand tangent cone of $K$ and
$\Pi_{T_K(x)}$ is the Euclidean projection onto it.

Define the residual function
\begin{equation}
\Psi(x)
:=
\mathrm{dist}^2\!\big(
0,\,
\partial_C \mathcal R(x)+N_K(x)
\big),
\label{eq:KKT_residual}
\end{equation}
where $N_K(x)$ denotes the Clarke normal cone to $K$.

Then the dissipation inequality
\begin{equation}
\sup_{v\in F(x)}
\mathcal R^\circ(x;v)
\le
-\Psi(x)
\qquad \text{for all } x\in K
\label{eq:dissipation_result}
\end{equation}
holds.
Consequently Assumption~\ref{ass:dissipation}
is satisfied with $\alpha=1$.
\end{proposition}

\begin{proof}
Fix $x\in K$.
Because $K$ is closed and convex,
$T_K(x)$ is a closed convex cone and its polar cone is
\[
N_K(x)=T_K(x)^\circ .
\]

Let $\zeta\in\partial_C\mathcal R(x)$ and define
\[
v := \Pi_{T_K(x)}(-\zeta).
\]
Since $T_K(x)$ is closed and convex,
the Euclidean projection is uniquely defined. By Moreau's decomposition theorem for convex cones,
there exists a unique $\eta\in N_K(x)$ such that
\begin{equation}
-\zeta = v+\eta,
\qquad
\langle v,\eta\rangle = 0.
\label{eq:Moreau_decomposition}
\end{equation}

Rearranging yields
\[
v=-\zeta-\eta,
\qquad
\zeta+\eta=-v.
\]

Because $\zeta\in\partial_C\mathcal R(x)$ and $\eta\in N_K(x)$,
we obtain
\[
\zeta+\eta \in \partial_C\mathcal R(x)+N_K(x).
\]

Moreover, by orthogonality in
\eqref{eq:Moreau_decomposition},
\begin{equation}
\|\zeta+\eta\|^2=\|v\|^2 .
\label{eq:projection_norm}
\end{equation}

Now compute the Clarke directional derivative.
Since $\mathcal R^\circ(x;\cdot)$ is convex and positively homogeneous,
for any $v\in F(x)$ there exists $\zeta\in\partial_C\mathcal R(x)$
such that
\[
\mathcal R^\circ(x;v)
=
\max_{\xi\in\partial_C\mathcal R(x)}
\langle \xi,v\rangle
\le
\langle \zeta,v\rangle .
\]

Using $v=-\zeta-\eta$ gives
\[
\langle\zeta,v\rangle
=
-\|\zeta\|^2-\langle\zeta,\eta\rangle .
\]

Taking the inner product of
\eqref{eq:Moreau_decomposition}
with $\eta$ yields
\[
0
=
\langle v,\eta\rangle
=
\langle -\zeta-\eta,\eta\rangle
=
-\langle\zeta,\eta\rangle-\|\eta\|^2 .
\]

Hence
\[
\langle\zeta,\eta\rangle=-\|\eta\|^2 .
\]

Substituting gives
\[
\langle\zeta,v\rangle
=
-\|\zeta\|^2+\|\eta\|^2
=
-\|\zeta+\eta\|^2 .
\]

Using \eqref{eq:projection_norm} yields
\begin{equation}
\mathcal R^\circ(x;v)\le -\|v\|^2 .
\label{eq:directional_decay}
\end{equation}

By definition of $\Psi(x)$ in \eqref{eq:KKT_residual},
\[
\|v\|^2
=
\mathrm{dist}^2\!\big(
0,\,
\partial_C\mathcal R(x)+N_K(x)
\big)
=
\Psi(x).
\]

Combining this with \eqref{eq:directional_decay}
yields
\[
\mathcal R^\circ(x;v)
\le
-\Psi(x).
\]

Taking the supremum over $v\in F(x)$ gives
\eqref{eq:dissipation_result}.
\end{proof}

\end{document}
